\section{Rational minimality and fixed periods}\label{sec:rational}\label{altord:h:section}
Our objective is to recognize when a symmetry that fixes all highest
lines must fix the full periods. We first discard the rational
subvariations that do not contribute to the highest line. Rationality
is essential: the invariant spaces of rational normal monodromy
subgroups will then be Hodge substructures.

Let $U$ be a smooth connected complex algebraic variety and let
$V=(V_{\Z},F^\bullet V_{\mathcal O},Q)$ be a polarized integral pure
variation of Hodge structure of weight $w$. Integral means that the
underlying rational local system has a monodromy-invariant lattice;
the polarization is rational and need not be integral. Write $m$ for the largest index with
$F^mV_{\mathcal O}\ne0$, and assume throughout that
$\rk F^mV_{\mathcal O}=1$. We call this piece the \emph{Hodge line}.
On a simply connected flat chart its inclusion in $V_{\mathcal O}$
defines a holomorphic map
\[
 \ell:U_{\mathrm{loc}}\longrightarrow \mathbb P(V_{\C,u}).
\]
Changes of flat chart compose $\ell$ with a constant projective linear
transformation, so its differential rank is intrinsic.

\subsection{The smallest rational subvariation}

We call a lifted point \emph{Hodge-generic} if it lies outside every
proper Hodge locus of a rational flat tensor in tensor constructions on
the variation and its dual. These are countably many closed analytic
loci on the universal cover. Their complement is dense by the Baire
theorem. Only this analytic genericity is needed in the Hodge arguments.

\begin{lemma}[Rationally minimal variation]\label{altord:h:minimal}
There is a unique smallest rational subvariation $V^{\min}\subset V_{\Q}$
whose Hodge filtration contains $F^mV_{\mathcal O}$. At a Hodge-generic
point, its fibre is the smallest rational Hodge substructure containing
the Hodge line, and that fibre is simple. The subvariation has an induced
polarization and integral lattice. On a smooth compactification with
unipotent boundary monodromy, its extended Hodge line agrees with that
of $V$.
\end{lemma}

\begin{proof}
In a polarizable rational Hodge structure, intersect all rational Hodge
substructures containing the specified line. A finite intersection
already attains the minimum dimension, and gives the unique smallest
one, say $W$. Semisimplicity shows that $W$ is simple: in a decomposition
into simple summands exactly one summand can have a nonzero deepest
Hodge piece, since that piece has dimension one, and that summand already
contains the line.

Here is also the generic-fibre justification, which is useful when
passing between variations and individual Hodge structures. Trivialize
the rational local system on the connected universal cover of $U$.
For each rational endomorphism $e$, the condition that $e$ preserve every
$F^p$ is closed analytic. Choose a lifted point outside all such loci
which are proper. A rational Hodge projector onto $W$ at this point
therefore preserves the filtration everywhere on the cover. Its action
on the Hodge line is the identity everywhere: it is an idempotent on a
line bundle, and equals the identity at the chosen point. Its constant
image is thus a subvariation containing the line. At every other point
chosen in the same way, applying the argument in both directions shows
that its image is again the smallest Hodge substructure containing the
line. Uniqueness and the deck action now imply monodromy invariance,
so this subvariation descends to $U$. Any rational subvariation containing
the line contains it, by inspection at the chosen point.

Restriction of $Q$ to a Hodge substructure is nondegenerate, and
$V^{\min}\cap V_{\Z}$ supplies the lattice. The polarization gives a
complementary subvariation whose deepest Hodge piece is zero.
Canonical extensions preserve this direct sum, proving the last claim.
This is the pure case of the transcendental-part construction of
\cite[Definition~2.9]{BFMT}.
\end{proof}

\subsection{From a fixed highest line to a fixed full period}
A variation is \emph{minimal for its highest line} if the generic
fibre is the smallest rational Hodge substructure containing that
line. Lemma~\ref{altord:h:minimal} supplies this replacement without
changing the line. The next two lemmas explain what this minimality
buys. We use the usual Mumford--Tate group of a rational Hodge
structure: the smallest rational algebraic group through which its
Hodge homomorphism factors.

\begin{lemma}[Normal subgroups and the deepest line]
\label{altord:p:normality}
Let $\mathbb V$ be a rationally polarized integral pure variation on a smooth
connected complex algebraic variety, minimal for its deepest Hodge line.
In a flat trivialization on the universal cover, let $G\subset\mathrm{GL}(V)$
be connected algebraic monodromy and let $P=\operatorname{MT}(V)$ be the generic
Mumford--Tate group.  If $H\subset G$ is a connected normal algebraic
subgroup defined over $\Q$, and $V^H_{\C}$ contains the deepest line at a
Hodge-generic lifted point, then $H$ is trivial.
\end{lemma}

\begin{proof}
The monodromy and fixed-part theorems give that $G$ is semisimple and
normal in the connected reductive group $P$; see
\cite[\S5, Theorem~1 and Corollary~1]{AndreMonodromy}.  Over an
algebraic closure, a connected normal subgroup of a semisimple group is
the product of a collection of its almost simple factors.  Conjugation
by $P$ preserves $G$.  Since $P$ is connected, it acts trivially on the
finite set of these factors, and therefore preserves $H$.

Consequently $V^H$ is a rational $P$-subrepresentation.  It is a rational
Hodge substructure at a Hodge-generic point, and it contains the deepest
Hodge line.  Minimality gives $V^H=V$.  The inclusion
$G\subset\mathrm{GL}(V)$ is faithful, so $H=1$.
\end{proof}

\begin{lemma}[A flat automorphism over a fixed point]
\label{altord:p:inertia}
Let $U$ be a smooth connected complex algebraic variety with a rationally
polarized integral pure variation $\mathbb V$ minimal for its deepest Hodge line.
Suppose that an automorphism $b$ of $U$ fixes a point $e$, and that there
is a flat isomorphism of variations
\[
  \Phi_y:V_y\longrightarrow V_{b(y)}.
\]
Use the fibre $V_e$ as the flat reference space, and put $a=\Phi_e$.
If, on a nonempty open subset of the universal cover, all deepest Hodge
lines lie in a single eigenspace of $a$, then $a$ preserves the complete
Hodge filtration at every lifted point.  In particular, locally at $e$
the full lifted period maps at $y$ and $b(y)$ are equal.
\end{lemma}

\begin{proof}
Write $\rho$ for monodromy based at $e$.  Flatness and $b(e)=e$ give
\[
  a\rho(\gamma)a^{-1}=\rho(b_*\gamma)
  \qquad (\gamma\in\pi_1(U,e)).
\]
Thus $a$ normalizes both algebraic monodromy and its identity component
$G$.  Let $\lambda$ be the eigenvalue in the hypothesis.  The condition
that the deepest line lies in $\ker(a-\lambda)$ is a holomorphic linear
condition on the universal cover, so it holds throughout that connected
cover by analytic continuation.

Fix a Hodge-generic lifted point and let $W\subset V_{\C}$ be the complex span
of the monodromy orbit of its deepest line.  This space is monodromy
invariant, hence $G$-invariant. Equivariance places every line in this
orbit among the deepest lines on the universal cover, so
$a|_W=\lambda\,\mathrm{id}_W$.
Every element $aga^{-1}g^{-1}$, with $g\in G$, therefore acts trivially
on $W$, as do all its $G$-conjugates.  Let $H$ be their normal algebraic
closure in $G$.  It is defined over $\Q$, since $a$ and $G$ are rational.
It is connected: the commutator map from the connected group $G$ has
connected image containing the identity, and the algebraic subgroup
generated by its conjugates is the closure of products of such
connected sets. Although $W$ need not be rational, $V^H$ is rational
and contains the chosen deepest line.
Lemma~\ref{altord:p:normality} gives $H=1$.  Hence $a$ centralizes $G$.

Pass to the connected finite \'etale cover $U_1\to U$ corresponding to
the inverse image of $G$ under monodromy, and choose $e_1$ over $e$.
Since $a$ centralizes all monodromy on $U_1$, it extends to a global flat
rational section of $\End(\mathbb V|_{U_1})$.  The latter is a polarizable
pure variation of weight zero. The flat section $a$ has Hodge type
$(0,0)$ at $e_1$, since $\Phi_e$ is a Hodge isomorphism. The base
$U_1$ is a smooth algebraic variety and hence a Zariski open subset
of a compact complex variety, so the fixed-part theorem applies.
By \cite[Theorem~7.22 and Corollary~7.23]{Schmid}, $a$ has type
$(0,0)$ everywhere. Thus $a$ preserves the full filtration on the
universal cover.

For completeness, take the lift of $b$ to the universal cover fixing a
chosen lift of $e$.  Flatness of $\Phi$ gives
$F^\bullet(b\widetilde y)=aF^\bullet(\widetilde y)$ in the reference
space.  The assertion just proved makes the right-hand side equal to
$F^\bullet(\widetilde y)$.
\end{proof}
